\section{SQL}
\begingroup
\setminted{fontsize=\fontsize{6.5}{7}\selectfont,
  breaklines=true, baselinestretch=1}

\subsection*{DDL}

\begin{minted}[style=manni]{sql}
--## create a table ##--
CREATE TABLE TableName (
    col1 datatype1 UNIQUE,
    col2 datatype2 NOT NULL,
    ...,
    PRIMARY KEY (col1, ...), -- PK -> NOT NULL
    FOREIGN KEY (col2, ...)
        REFERENCES OtherTable(col1, ...)
);

--## common datatypes ##--
-- numeric: INT, SMALLINT, BIGINT, DECIMAL(p,s),
--          NUMERIC(p,s), REAL, DOUBLE PRECISION
-- character: CHAR(n), VARCHAR(n), TEXT
-- date/time: DATE, TIME, TIMESTAMP

--## delete a table ##--
DROP TABLE TableName; -- dependent FKs may block

--## modify a table ##--
ALTER TABLE TableName ADD newCol datatype;
ALTER TABLE TableName DROP colToDelete;
ALTER TABLE TableName ADD PRIMARY KEY (col1, ..);
ALTER TABLE TableName ADD FOREIGN KEY (col2, ..)
    REFERENCES OtherTable(col1, ..);
\end{minted}

\subsection*{DML}

\begin{minted}[style=manni]{sql}
--## insert a row ##--
INSERT INTO TableName VALUES
    (val1, val2, ...), -- one row
    (val1, val2, ...),
    ...;

--## load from a file ##--
LOAD DATA LOCAL INFILE 'file.txt'
INTO TABLE TableName
FIELDS TERMINATED BY ','
LINES TERMINATED BY '\n';

--## delete row(s) ##--
DELETE FROM TableName; -- delete all rows
DELETE FROM TableName  -- delete rows conditionally
WHERE condition; 

--## update row(s) ##--
UPDATE TableName SET col1 = val1, col2 = val2, ...
WHERE condition; -- update rows conditionally
UPDATE TableName SET col1 = col1 * 2; -- arithmetic

--## case statement ##--
UPDATE TableName SET col1 = CASE
    WHEN condition1 THEN val1
    WHEN condition2 THEN val2
    ELSE val3
END; -- CASE can be used elsewhere, e.g. SELECT

--## query rows ##--
SELECT col1, col2, ... FROM TableName;
SELECT * FROM TableName;
SELECT * FROM TableName WHERE condition;

--## string matching ##--
SELECT * FROM T WHERE col LIKE 'hello';
    -- matches 'hello' exactly
SELECT * FROM T WHERE col LIKE 'hello%';
    -- matches 'hello', 'hello world', ...
SELECT * FROM T WHERE col LIKE 'h_llo';
    -- matches 'hello', 'hallo', 'hxllo', ...
\end{minted}
\subsection*{Query}
\begin{minted}[style=manni]{sql}
--## select / alias / group / sort ##--
SELECT DISTINCT dept, COUNT(*) AS n
FROM Employee E
WHERE salary > 0       -- filter rows
GROUP BY dept
HAVING COUNT(*) >= 2   -- filter groups
ORDER BY n DESC, dept ASC; -- default ASC
-- order: FROM -> WHERE -> GROUP BY -> HAVING
--        -> SELECT/DISTINCT -> ORDER BY

--## aggregate ##--
SELECT SUM(x), AVG(x), MIN(x), MAX(x),
       COUNT(*), COUNT(x), COUNT(DISTINCT x)
FROM T;
-- COUNT(*): rows; others: ignore NULL
-- empty input: COUNT -> 0; others -> NULL
SELECT SUM(CASE WHEN x > 0 THEN x ELSE 0 END)
FROM T; -- conditional sum

--## join ##--
SELECT A.*, B.* FROM A
INNER JOIN B ON A.id = B.aid; -- matched pairs
-- LEFT/RIGHT JOIN: keep all left/right rows
-- unmatched side -> NULL; CROSS JOIN: m*n rows
SELECT A.id, COUNT(B.id) FROM A
LEFT JOIN B ON A.id = B.aid
GROUP BY A.id; -- no matches -> 0, not COUNT(*)

--## set operations ##--
SELECT id FROM A UNION SELECT id FROM B;
-- UNION: either; INTERSECT: both; EXCEPT: A-B
-- same #columns + compatible types
-- duplicates removed; UNION ALL keeps them
\end{minted}

\subsection*{Nested Query}
\begin{minted}[style=manni]{sql}
--## membership ##--
SELECT * FROM Employee
WHERE dept IN (SELECT id FROM Dept
               WHERE budget > 100);
-- IN / NOT IN: member / not member

--## scalar: 1 column, <=1 row ##--
SELECT * FROM Employee
WHERE salary > (SELECT AVG(salary) FROM Employee);
-- 0 rows -> NULL; >1 row -> error

--## correlated: uses outer alias ##--
SELECT D.id, (SELECT COUNT(*) FROM Employee E
              WHERE E.dept = D.id) AS n
FROM Dept D; -- no matches -> 0

--## exists / not exists ##--
SELECT D.* FROM Dept D
WHERE EXISTS (SELECT 1 FROM Employee E
              WHERE E.dept = D.id);
-- EXISTS: >=1 row; NOT EXISTS: 0 rows
-- EXISTS(SELECT COUNT(*) ...) -> true
-- (without GROUP BY/HAVING)

--## all / some (any) ##--
SELECT * FROM Employee
WHERE salary >= ALL (SELECT salary FROM Employee);
-- ALL: every; SOME/ANY: at least one
-- empty set: ALL -> true; ANY -> false
-- = ANY: IN; <> ALL: NOT IN

--## maximum: keep ties ##--
SELECT * FROM Employee
WHERE salary = (SELECT MAX(salary) FROM Employee);
-- MAX ignores NULL; ALL does not
-- choose population: global vs eligible only

--## derived table: max of group counts ##--
SELECT MAX(T.n) FROM
    (SELECT dept, COUNT(*) AS n FROM Employee
     GROUP BY dept) AS T;
\end{minted}

\subsection*{NULL \& UNKNOWN}
\begin{minted}[style=manni]{sql}
SELECT * FROM T WHERE x IS NULL; -- not = NULL
SELECT * FROM T WHERE x IS NOT NULL;
-- arithmetic with NULL -> NULL
-- comparison with NULL -> UNKNOWN (U)
-- WHERE/HAVING: only TRUE survives
--     OR      AND   NOT
--   T U F    T U F   
-- T|T T T  T|T U F  T|F
-- U|T U U  U|U U F  U|U
-- F|T U F  F|F F F  F|T
\end{minted}

\subsection*{Useful Patterns}
\begin{minted}[style=manni]{sql}
--## every child: x > 0 (NULL fails) ##--
SELECT P.* FROM Parent P
WHERE NOT EXISTS (SELECT 1 FROM Child C
    WHERE C.pid = P.id
      AND (C.x <= 0 OR C.x IS NULL));
-- empty set passes; add EXISTS if nonempty

--## every required item has a match ##--
SELECT P.* FROM Person P
WHERE NOT EXISTS (SELECT 1 FROM Required R
    WHERE NOT EXISTS (SELECT 1 FROM Match M
        WHERE M.pid = P.id AND M.rid = R.id));

--## exactly k distinct items ##--
SELECT pid FROM Match GROUP BY pid
HAVING COUNT(DISTINCT rid) = 2;
-- joins may duplicate rows: count IDs, not rows
\end{minted}

\subsection*{View / Authorization / Assertion}
\begin{minted}[style=manni]{sql}
--## view ##--
CREATE VIEW V AS SELECT id, name FROM Employee;

--## privileges / roles ##--
GRANT SELECT ON TableName TO username;
GRANT UPDATE(col) ON TableName TO roleName;
REVOKE SELECT ON TableName FROM username;
CREATE ROLE roleName;
GRANT roleName TO username;

--## assertion (standard SQL; not MySQL) ##--
CREATE ASSERTION AssertionName CHECK (
    (SELECT COUNT(*) FROM Dept) <=
    (SELECT COUNT(*) FROM Employee)
); -- DB-wide condition; check on updates
\end{minted}
\endgroup
